% LTeX: enabled=false

% ---------------------------------------------------------------
% Estilo (plantilla de las presentaciones de IdS2 y PDS)
% ---------------------------------------------------------------
\usefonttheme{serif}
\definecolor{azul}{HTML}{2C6AA0}
\definecolor{verde}{HTML}{3E8E5E}
\definecolor{gris}{HTML}{F1F1F1}
\definecolor{rojo}{HTML}{B03A2E}

\setbeamercolor{frametitle}{fg=azul}
\setbeamercolor{title}{fg=azul}
\setbeamercolor{structure}{fg=azul}
\setbeamercolor{alerted text}{fg=azul}
\setbeamerfont{alerted text}{series=\bfseries}
\setbeamercolor{block title}{bg=azul, fg=white}
\setbeamercolor{block body}{bg=gris}
\setbeamerfont{frametitle}{series=\bfseries, size=\Large}
\setbeamertemplate{navigation symbols}{}
\setbeamertemplate{itemize items}{\color{azul}\textbullet}
\setbeamertemplate{blocks}[default]

\setlength{\leftmargini}{1.2em}   % listas de primer nivel
\setlength{\leftmarginii}{1.2em}  % listas anidadas

% Pie de página: autores, título corto y número de diapositiva.
% Tres cajas que suman exactamente el ancho de la hoja (0.35 + 0.40 + 0.25).
% Cada línea termina en % para que los saltos de línea no agreguen espacios.
\setbeamercolor{footline}{fg=gray}
\setbeamerfont{footline}{size=\scriptsize}
\setbeamertemplate{footline}{%
    \leavevmode%
    \hbox{%
        \begin{beamercolorbox}[wd=0.35\paperwidth, ht=3ex, dp=1.2ex, left]{footline}%
            \hspace*{1em}Aranda - Wlodeck%
        \end{beamercolorbox}%
        \begin{beamercolorbox}[wd=0.40\paperwidth, ht=3ex, dp=1.2ex, center]{footline}%
            KERALYTIX - Proyecto final integrador%
        \end{beamercolorbox}%
        \begin{beamercolorbox}[wd=0.25\paperwidth, ht=3ex, dp=1.2ex, right]{footline}%
            \insertframenumber\,/\,\inserttotalframenumber\hspace*{1em}%
        \end{beamercolorbox}%
    }%
    \vskip0pt%
}

% Título de cada diapositiva con una línea debajo
\setbeamertemplate{frametitle}{%
    \vspace{0.4em}\insertframetitle\par\vspace{0.2em}{\color{azul}\hrule height 0.6pt}}

% ---------------------------------------------------------------
% Macros propias
% ---------------------------------------------------------------

% Idea clave: caja recta, fondo gris y borde azul
\newcommand{\idea}[1]{%
    \begin{center}
        \setlength{\fboxsep}{5pt}%
        \fcolorbox{azul}{gris}{\parbox{0.9\linewidth}{\centering #1}}
    \end{center}}

% Marca de contenido pendiente (no debe quedar ninguna en la versión final)
\newcommand{\pendiente}[1]{%
    \fcolorbox{rojo}{rojo!5}{\parbox{0.9\linewidth}{\centering\small\color{rojo}%
    \textbf{PENDIENTE}\\ #1}}}

% Estado de un resultado, como pide la consigna: alcanzado, en proceso o pendiente.
% #1 = color, #2 = estado, #3 = objetivos que cubre. Va al final del frame.
\newcommand{\estado}[3]{%
    \par\vfill\hfill{\small\fcolorbox{#1}{#1!10}{\textcolor{#1}{\textbf{#2}}}\enspace\textcolor{gray}{#3}}}
\newcommand{\alcanzado}[1]{\estado{verde}{Alcanzado}{#1}}
\newcommand{\enproceso}[1]{\estado{azul}{En proceso}{#1}}
\newcommand{\porhacer}[1]{\estado{gray}{Por empezar}{#1}}

% Etiqueta de estado dentro de una tabla
\newcommand{\etiq}[2]{\fcolorbox{#1}{#1!10}{\textcolor{#1}{\textbf{#2}}}}

% Formato común de las tablas: encabezado azul y líneas horizontales grises
\NewTblrEnviron{tabla}
\SetTblrOuter[tabla]{halign=c}
\SetTblrInner[tabla]{row{1} = {bg=azul, fg=white, font=\bfseries},
                     hlines = {gray!50}, rowsep = 2pt}

% Sección: solo ordena el archivo, sin diapositiva separadora,
% para no gastar páginas (límite de 20 páginas y 20 minutos)
\newcommand{\seccion}[1]{\section{#1}}

% Estilos de los diagramas TikZ
% Estilos comunes de los diagramas (ver Plantilla_beamer.md, sección 4)
\tikzset{
    caja/.style={draw=azul, fill=white, align=center, font=\small,
                 execute at begin node={\hyphenpenalty=10000},
                 inner sep=5pt, text width=#1},
    caja/.default=3cm,
    flecha/.style={-{Stealth[length=2.5mm]}, thick, azul},
}

